% Classical low-degree standard-D cases with a complete derivation.
\paragraph{C.2k (Known with derivation: the exact low-degree standard-$D$ range used here).}
Let $X/\mathbb C$ be smooth projective of dimension $n$ and let
$A^r_H(X)\subset H^{2r}(X,\mathbb Q)$ be the space of algebraic
cycle classes. Numerical and homological equivalence, with rational
coefficients, agree in every codimension
\[
                        r\in\{0,1,2,n-1,n\}.
\]
Repeated or out-of-range entries are omitted. In particular the
full standard conjecture $D(X)$ holds for every complex variety
of dimension at most four. This is not an assertion of the
Lefschetz standard conjecture for all correspondences on those
varieties.

\begin{proof}
We derive these classical cases directly from the opened
Lefschetz $(1,1)$ and Hodge--Riemann inputs
\citep[Remark 3.4, p.~379, and Remark 3.10, p.~380]{cyKleiman1968}.
The degree map treats $r=0,n$. For divisors, the Hodge index
form
\[
       (d,e)\longmapsto\int_X deh^{n-2}
\]
is nondegenerate on the rational Neron--Severi space; for curves
of dimension one we instead use degree directly. Thus a
numerically trivial divisor has zero cohomology class.

For $r=n-1$ and $n\ge2$, cohomological hard Lefschetz is an
isomorphism of rational Hodge structures
\[
 h^{n-2}:H^2(X,\mathbb Q)\xrightarrow{\sim}
                         H^{2n-2}(X,\mathbb Q).
\]
The inverse of an algebraic class is a rational $(1,1)$ class,
and therefore is a divisor class. We have
$A_H^{n-1}(X)=h^{n-2}A_H^1(X)$. Pairing with divisors and
the preceding nondegenerate form proves the assertion.

It remains to treat $r=2$ with $n\ge4$.
Take $a\in A_H^2(X)$ which is numerically trivial.
There is a unique rational class $d\in H^2(X,\mathbb Q)$ such
that
\[
                      h^{n-2}d=h^{n-3}a.
\]
Hard Lefschetz and Hodge type make $d$ a rational $(1,1)$
class, hence an algebraic divisor class. For every divisor
class $e$ we have
\[
             \int_X deh^{n-2}=
             \int_X aeh^{n-3}=0.
\]
The divisor Hodge index form gives $d=0$. Thus
$h^{n-3}a=0$: the class $a$ is primitive in degree four.
Numerical triviality also gives
\[
                         \int_X a^2h^{n-4}=0.
\]
The primitive Hodge--Riemann form is positive definite on
rational $(2,2)$ classes, so $a=0$. These arguments prove
the stated range. For $n\le4$ that range includes every
codimension.
\end{proof}

There is a precise remaining quotient in degree two of the
homological Lefschetz problem. For every $n\ge4$, the pairing
\[
       Q_2(a,b)=\int_X abh^{n-4}
                  \quad(a,b\in A_H^2(X))
\]
is nondegenerate. To verify this without assuming a standard
conjecture, use the same rational algebraic projection
$a\mapsto d$ above. It gives the orthogonal decomposition
\[
 A_H^2(X)=
 \bigl(A_H^2(X)\cap\ker h^{n-3}\bigr)\oplus hA_H^1(X).
\]
The first summand has the positive Hodge--Riemann form and
the second has the nondegenerate divisor form. Define
\[
 R_H^{n-2}(X)=
 \left\{z\in A_H^{n-2}(X):
       \int_X za=0\text{ for every }a\in A_H^2(X)\right\}.
\]
Then there is an exact direct-sum decomposition of rational
cycle-class spaces
\begin{equation}\label{eq:C2k-first-obstruction}
 A_H^{n-2}(X)=h^{n-4}A_H^2(X)\oplus R_H^{n-2}(X).
\end{equation}
Indeed $Q_2$ represents the functional $a\mapsto\int_X za$
by a unique $b\in A_H^2(X)$. Subtracting $h^{n-4}b$ gives the
second summand; the same form shows that the intersection
of the two summands is zero. This is a finite-dimensional
rational linear-algebra argument on algebraic cycle classes,
without assuming that every rational Hodge class is algebraic.

Consequently $A_H(X,2)$ is equivalent to
$R_H^{n-2}(X)=0$, which is exactly standard $D$ in
codimension $n-2$ on this $X$. On a fivefold this reads
\[
 A_H^3(X)/hA_H^2(X)\simeq R_H^3(X).
\]
Thus the first dimension and codimension not covered by
the preceding classical theorem are $(n,r)=(5,3)$.
No vanishing of this last space is proved here.

Finally, C.1n's fixed-target endpoint image criterion still
quantifies over \emph{all source varieties and all source
cycles}. Its necessity argument can insert a numerically
trivial cycle of any source dimension and codimension.
The theorem just proved therefore does not turn that
universal image criterion, even for a point or projective
space as target, into an unconditional assertion.
